\section{Introduction}
\label{sec:intro}

Every game has a design DNA: a characteristic pattern of what it emphasizes relative
to what it could emphasize, expressed across the full space of design axes. Two games
with identical average scores across 32 axes are not identical games --- they differ in
which axes they prioritize, and that pattern of relative emphasis is what designers,
publishers, and players respond to when they describe a game as ``very Knizia'' or
``quintessential Lacerda''. This paper introduces the design fingerprint as a formal
object and the within-game z-score as its correct computational representation.

\subsection{Why Raw Scores Obscure Design Character}

When game scores are analyzed without normalization, the dominant signal is overall
quality: high-rated games score above the corpus mean on almost every axis, and
low-rated games score below. The \emph{pattern} of relative emphasis --- which axes
a game elevates above its own average --- is masked by this quality signal.
Within-game z-score normalization removes the quality signal by subtracting each
game's own mean before analysis, leaving only the relative emphasis pattern.

\subsection{Fingerprints as Design Objects}

A design fingerprint is the vector of within-game z-scores across all 32 axes,
representing each game's characteristic deviation from its own average.
Fingerprints are quality-neutral by construction: a simple filler game and a
complex euro game can have similar fingerprints if they emphasize the same
relative mix of design axes.
